Every candidate needs a primary distance estimate, so a scan repeatedly
looks up code-dependent contributions. Large tables consume on-chip
storage or require global-memory reads. JHQ's coordinate structure lets
the GPU use smaller tables while retaining the same quantiser. Primary
ADC uses FP32 queries $q'$ and stored reconstructions \cite{pq,jhq}:
\begin{equation}
 T_m[z]=\|q'_m-\mathcal C_m[z]\|_2^2,\qquad
 \primarydist(q,x)=\sum_{m=0}^{M-1}T_m[z_m(x)].
 \label{eq:5-primary-lut}
\end{equation}
For $K=256$ and $D_s=8$, the analytical codebook contains all sign
combinations at one reconstruction level:
\begin{equation}
 \begin{aligned}
 a&=\sigma\Phi^{-1}(3/4)\simeq0.67449\sigma,\qquad
 \mathcal C_m=\{-a,+a\}^{8},\\
 \primarydist(q,x)&=\|q'\|_2^2+da^2
       -2a\sum_{j=0}^{d-1}s_j(x)q'_j.
 \end{aligned}
 \label{eq:5-sign-codebook}
\end{equation}
Here $\Phi$ is the standard normal cumulative distribution function;
each byte stores eight signs, higher bits first. Primary ranking is
therefore equivalent to ranking the query's inner product with the signs:
\[
 a>0\quad\Longrightarrow\quad
 \arg\min_x\primarydist(q,x)=\arg\max_x\sum_j s_j(x)q'_j.
\]

\subsection{Factorisation without changing the quantiser}

The high and low four bits select independent coordinate halves. Write
$z=16h+l$, with $h,l\in\{0,\ldots,15\}$ and
$\mathcal C_m[z]=(A_m[h],B_m[l])$. Define one table for each half:
\begin{equation}
 T_m^{\rm hi}[h]=\|q'_{m,\rm hi}-A_m[h]\|_2^2,\qquad
 T_m^{\rm lo}[l]=\|q'_{m,\rm lo}-B_m[l]\|_2^2.
 \label{eq:5-lut-halves}
\end{equation}
The squared distance separates over these halves:
\begin{equation}
 T_m[z]=T_m^{\rm hi}[z\!\gg\!4]
        +T_m^{\rm lo}[z\mathbin{\&}15].
 \label{eq:5-lut-factorisation}
\end{equation}
\par\medskip\noindent\begin{minipage}{\linewidth}
\captionsetup{hypcap=false}
\centering
% Schematic indexing paths; exact identity is conditional on Cartesian coordinates.
\begin{tikzpicture}[x=1cm,y=1cm,font=\small\linespread{1}\selectfont,
 box/.style={draw=ThesisNavy!60,rounded corners=2pt,fill=ThesisNavy!4,inner sep=4pt,align=center},
 arr/.style={-{Stealth[length=4pt]},semithick,ThesisNavy!80}]
\node[font=\small] at (7.3,5.40)
 {Cartesian premise: $\mathcal C_m[z]=(A_m[h],B_m[l])$};
\node[box,fill=ThesisNavy!12] (byte) at (7.3,4.70) {$z=(hhhh\ llll)_2=16h+l$};
\node[anchor=west,font=\bfseries] at (.15,4.95) {Full table};
\node[anchor=east,font=\bfseries] at (14.7,4.95) {Two small tables};
\node[box] (fullindex) at (2.8,3.55) {index $z\in[0,255]$};
\node[box] (hiindex) at (9.25,3.55) {$h=z\gg4$\\$hhhh$: four signs\\$2^4=16$ patterns};
\node[box] (loindex) at (12.65,3.55) {$l=z\mathbin{\&}15$\\$llll$: four signs\\$2^4=16$ patterns};
\draw[arr] (byte.west)-|(fullindex.north);
\draw[arr] (byte.east)-|(loindex.north);
\draw[arr] (byte.south)-|(hiindex.north);

\draw[ThesisNavy!55] (1.05,1.40) rectangle (4.55,2.65);
\node at (2.8,2.42) {$T_m[0],\ldots,T_m[255]$};
\draw[ThesisNavy!35] (1.05,2.17)--(4.55,2.17);
\node[fill=ThesisNavy!12,minimum width=3.4cm,minimum height=.42cm]
 (fullread) at (2.8,1.87) {$T_m[z]$};
\draw[arr] (fullindex.south)--(2.8,2.65);
\node[align=center] at (2.8,.88) {256 entries\quad 1 read\\1,024 bytes (FP32)};

\foreach \xx/\name/\idx in {9.25/hi/h,12.65/lo/l} {
 \draw[ThesisNavy!55] (\xx-1.45,1.40) rectangle (\xx+1.45,2.65);
 \node at (\xx,2.42) {$T_m^{\rm \name}[0..15]$};
 \draw[ThesisNavy!35] (\xx-1.45,2.17)--(\xx+1.45,2.17);
 \node[fill=ThesisNavy!12,minimum width=2.8cm,minimum height=.42cm]
  (read\name) at (\xx,1.87) {$T_m^{\rm \name}[\idx]$};
}
\draw[arr] (hiindex.south)--(9.25,2.65);
\draw[arr] (loindex.south)--(12.65,2.65);
\node[draw=ThesisNavy,circle,inner sep=2pt] (sum) at (10.95,.88) {$+$};
\draw[arr] (9.25,1.40)|-(sum.west);
\draw[arr] (12.65,1.40)|-(sum.east);
\node[align=center] at (10.95,.10) {16+16 entries\quad 2 reads + 1 sum\\128 bytes (FP32): $1/8$ of full-table storage};
\draw[dashed,ThesisNavy!40] (6.25,.1)--(6.25,3.45);
\end{tikzpicture}

\captionof{figure}[Cartesian primary-LUT factorisation]{Cartesian primary-LUT factorisation.}
\label{fig:5-cartesian-lut}
\end{minipage}\par\medskip
Each byte now indexes two 16-entry tables instead of one 256-entry
table, reducing table storage eightfold. The identity holds in real
arithmetic; FP32 evaluation may associate additions differently. The
builder checks that codewords $16h,l$ select the intended coordinate
halves. Merely splitting the labels of a learned PQ codebook into
nibbles would not establish this property.

\subsection{Grouping, construction work and resource cost}

More generally, let $\glut\in\{1,2,4,8\}$ divide the eight coordinates
into contiguous groups whose bit fields match the Cartesian digits.
For $D_s=8$, each group has $2^{8/\glut}$ entries. The total entries
per subspace and FP32 bytes per query are
\begin{equation}
 n_{\rm LUT}(\glut)=\glut 2^{8/\glut},\qquad S_{\rm LUT}(\glut)=4Mn_{\rm LUT}(\glut)\ \text{bytes}.
 \label{eq:5-lut-cost}
\end{equation}
Constructing the entries independently evaluates $D_s2^{8/\glut}$
coordinate contributions per subspace. Two groups need one eighth of
the full-table storage and one sixteenth of its coordinate work.
\Cref{tab:5-lut-resources} also shows the cost of joining the lookups;
these algebraic counts do not predict whole-query speedups.

\par\medskip\noindent\begin{minipage}{\linewidth}
\captionsetup{hypcap=false}
\centering
\captionof{table}[Cartesian lookup-table trade-offs]{Cartesian table trade-offs for an eight-bit, eight-coordinate
primary subspace. Combining additions are the algebraic minimum for
joining the looked-up values, excluding accumulation across subspaces;
they are not measured instruction counts.}
\label{tab:5-lut-resources}
\small
\begin{tabular}{@{}rrrrrr@{}}
\toprule
$\glut$ & Entries & FP32 bytes & Lookups & Combining adds & Build terms \\
\midrule
1 & 256 & 1,024 & 1 & 0 & 2,048 \\
2 & 32 & 128 & 2 & 1 & 128 \\
4 & 16 & 64 & 4 & 3 & 32 \\
8 & 16 & 64 & 8 & 7 & 16 \\
\bottomrule
\end{tabular}
\end{minipage}\par\medskip
Placement depends on the table and selection buffer together. Tables
are built in global workspace and copied cooperatively into shared
memory when the complete allocation fits the per-block limit
$S_{\max}$. For buffer capacity $C_{\rm buf}$, the ordinary and grouped
scan allocations are
\begin{equation}
 S_0=8C_{\rm buf}+8+S_{\rm LUT}(\glut),\qquad
 S_{\rm grp}=8C_{\rm buf}+24+S_{\rm LUT}(\glut)
 \quad\text{bytes}.
 \label{eq:5-shared-footprint}
\end{equation}
For the evaluated paths, without optional fused-refinement or early-exit
scratch, the launch rule is
\[
 \text{shared lookup}\iff S_{\rm path}\leq S_{\max}.
\]
The sweep uses $c=1{,}000$, blocks of 512 or 1,024 threads, and
$C_{\rm buf}=2{,}048$. Full tables occupy 96--384\,KiB, whereas two-group
tables occupy 12--48\,KiB. NVIDIA documents a 99\,KiB per-block limit
for compute capability 12.0 \cite{blackwell_tuning}. Even the smallest
full table exceeds this limit once its selection buffer is included:
at $M=96$, the ordinary path needs $114{,}696$ bytes. The largest
two-group allocation is $65{,}560$ bytes in the grouped path. These
footprints place full tables in global memory and two-group tables in
shared memory under the launch rule.\footnote{This placement follows
from the source rule and published limit; the archived measurements
do not retain a device-attribute query.}

Four and eight groups each need 16 entries per subspace, half as many
as two groups, but require more lookups and additions. Smaller tables
can free resources for concurrent blocks; extracting more bit fields,
joining their values and staging a copy for each scan group add work.
The small-table idea has precedents in \cite{fastscan,quickeradc,ivfrabitq};
here it is integrated with JHQ's two-level GPU search. \Cref{sec:7-3}
compares throughput and candidate-level FP32 differences among the
table forms.
